% Einheit 2 – Extrempunkte
\setcounter{aufgabe}{13}
\hypertarget{e2}{}\einheitenkopf{Einheit 2 von 5 · Extrempunkte}
\zweigzeile{Hier lernst du, Hoch-, Tief- und Sattelpunkte zu berechnen und nachzuweisen · kennst du seit Q1 (Klasse 11) · Abitur GK · baut auf: Monotonie (Einheit 1), Ableitungen bilden, Gleichungen lösen, Funktionswerte berechnen, Graphen grob skizzieren}

\begin{aufgabe}{Ich erkenne, ob eine Stelle gegeben ist oder gesucht wird. Lies die Aufgabenstellung und kreuze an. Rechne nichts.}
\begin{teile}
\teil „Bestimmen Sie die Koordinaten der Extrempunkte des Graphen von $f$.“ \\ \kreuz{gegeben – nachweisen}\kreuz{gesucht – berechnen}
\teil „Zeigen Sie, dass der Graph von $f$ an der Stelle $x = 2$ einen Tiefpunkt hat.“ \\ \kreuz{gegeben – nachweisen}\kreuz{gesucht – berechnen}
\teil „Weisen Sie nach, dass $H(1 \mid 4)$ ein Hochpunkt des Graphen von $f$ ist.“ \\ \kreuz{gegeben – nachweisen}\kreuz{gesucht – berechnen}
\teil „Ermitteln Sie die Stelle, an der der Graph von $f$ seinen Hochpunkt hat.“ \\ \kreuz{gegeben – nachweisen}\kreuz{gesucht – berechnen}
\teil „Untersuchen Sie, ob der Graph von $f$ bei $x = -1$ eine waagerechte Tangente hat.“ \\ \kreuz{gegeben – nachweisen}\kreuz{gesucht – berechnen}
\end{teile}
\end{aufgabe}

\begin{aufgabe}{Ich erkenne, welcher Nachweis für die Art eines Extrempunkts passt. Schreibe A (Vorzeichen von $f''$), B (Vorzeichenwechsel von $f'$), C (Begründung aus Abbildung oder Sachzusammenhang) oder D (kein Nachweis verlangt). Rechne nichts.}
\begin{teile}
\teil An der Stelle $x_0$ gilt $f'(x_0) = 0$ und $f''(x_0) = -4$. \leerfeld
\teil An der Stelle $x_0$ gilt $f'(x_0) = 0$ und $f''(x_0) = 0$. \leerfeld
\teil „Bestimmen Sie die Stelle des Hochpunkts rechnerisch. Die hinreichende Bedingung muss nicht untersucht werden.“ \leerfeld
\teil „Berechnen Sie die Stelle, an der die Konzentration am größten ist. Begründen Sie die Art des Extremums mithilfe der Abbildung.“ \leerfeld
\teil Gegeben ist nur $f'(x) = (2 - x)\cdot\mathrm{e}^{-x}$, die zweite Ableitung ist nicht angegeben. \leerfeld
\end{teile}
\end{aufgabe}

\begin{aufgabe}{Ich kann die Stellen berechnen, an denen ein Graph eine waagerechte Tangente hat. Bilde $f'$ und löse $f'(x) = 0$.}
\beispiel{f(x) &= x^3 - 6x^2 + 9x \\ f'(x) &= 3x^2 - 12x + 9 = 0 &&\mid :3 \\ x^2 - 4x + 3 &= 0 \\ x &= 2 \pm \sqrt{4 - 3} \\ x_1 = 3,\ x_2 &= 1}
\begin{gleichungsraster}[3]
\gl{f(x) = \tfrac{1}{3}x^3 - 4x} & \gl{f(x) = \tfrac{1}{3}x^3 - x^2 - 3x} \\
\gl{f(x) = 2x^3 + 3x^2 - 12x} & \gl{f(x) = -\tfrac{1}{3}x^3 + 2x^2 + 5x} \\
\end{gleichungsraster}
\end{aufgabe}

\begin{aufgabe}{Ich kann Hoch-, Tief- und Sattelpunkte berechnen. Bestimme Lage und Art der Extrempunkte.}
\begin{teile}
\teil $f(x) = x^3 + 3x^2 - 9x$
\teil $f(x) = x^4 - 4x^3 + 4x^2$
\teil $f(x) = 3x^4 - 4x^3$; untersuche auch, ob ein Sattelpunkt vorliegt.
\teil $f(x) = 0{,}25x^4 - 3{,}5x^2 + 6x$; die Stelle $x = 1$ ist eine Extremstelle. Bestimme alle Extremstellen und ihre Art.
\teil $f(x) = (x^2 - 3)\cdot\mathrm{e}^{x}$
\teil Die Zahl der Anrufe pro Minute bei einer Hotline wird durch $a(t) = 5t\cdot\mathrm{e}^{-0{,}2t}$ beschrieben ($t$ in Minuten nach Beginn einer Werbesendung). Im Sachzusammenhang ist klar, dass die Zahl der Anrufe erst steigt und dann wieder fällt; die hinreichende Bedingung muss nicht untersucht werden. Berechne, wann die meisten Anrufe pro Minute eingehen, und gib ihre Zahl an.
\teil Bestimme den größten und den kleinsten Funktionswert von $g(x) = -x^3 + 6x^2$ auf dem Intervall $[-3;\,5]$.
\teil Bestimme den Wertebereich von $f(x) = x^4 - 2x^2 + 3$.
\teil Bestimme Lage und Art aller lokalen Extrempunkte des Graphen von $f$ mit $f(x) = x^2\cdot\mathrm{e}^{-0{,}5x}$. (Abitur 2021)
\end{teile}
\end{aufgabe}

\begin{aufgabe}{Ich kann nachweisen, dass an einer gegebenen Stelle ein Hoch- oder Tiefpunkt liegt. Zeige, dass $f'(x_0) = 0$ gilt.}
\beispiel{f(x) &= x^3 - 6x^2 + 5, \quad x_0 = 4 \\ f'(x) &= 3x^2 - 12x \\ f'(4) &= 3\cdot 16 - 12\cdot 4 = 0 && \text{notwendige Bedingung erfüllt}}
\begin{teilezwei}
\tz $f(x) = x^2 - 10x$, \quad $x_0 = 5$ & \tz $f(x) = x^3 - 27x$, \quad $x_0 = 3$ \\
\tz $f(x) = 2x^3 + 3x^2 - 36x$, \quad $x_0 = 2$ & \tz $f(x) = x^4 - 18x^2$, \quad $x_0 = -3$ \\
\end{teilezwei}
\begin{teile}
\teil Zeige, dass der Graph von $f(x) = -x^3 + 3x^2 + 24x$ bei $x = 4$ einen Hochpunkt hat.
\teil Weise nach, dass $T(-1 \mid -7)$ ein Tiefpunkt des Graphen von $f(x) = x^4 + 4x - 4$ ist.
\teil Für $f(x) = (x + 2)^4 - 1$ ist $f''(-2) = 0$. Zeige mit dem Vorzeichenwechsel von $f'$, dass bei $x = -2$ ein Tiefpunkt liegt.
\teil Für eine Funktion $f$ gilt $f'(3) = 0$ und $f''(3) = -2$. Begründe, dass der Graph von $f$ bei $x = 3$ einen Hochpunkt hat.
\teil Zeige, dass der Graph von $f(x) = -x^3 + 3x^2 - 3x + 4$ bei $x = 1$ eine waagerechte Tangente, aber keinen Extrempunkt hat.
\teil Zeige, ohne $f'(x) = 0$ zu lösen, dass der Graph von $f(x) = x^3 + x^2 - 3x$ zwischen $x = 0$ und $x = 1$ einen Tiefpunkt hat.
\teil Gegeben ist $f(x) = 3x^4 - 8x^3 + 2$. Zeige, dass der Graph an der Stelle $x = 2$ einen Tiefpunkt hat, und gib ihn an. Begründe, dass der Graph an der Stelle $x = 0$ eine waagerechte Tangente, aber keinen Extrempunkt hat. (Abitur 2026)
\end{teile}
\end{aufgabe}

\begin{aufgabe}{Ich kann mit berechneten Hoch- und Tiefpunkten weiterrechnen.}
\begin{teile}
\teil Berechne die Nullstellen und die Extremstelle von $p(x) = x^2 - 6x + 5$.
\teil Bestimme den kleinsten senkrechten Abstand zwischen dem Graphen von $q(x) = x^2 - 4x + 7$ und der Geraden $y = 1$. Beschreibe dein Vorgehen.
\teil Weise nach, dass $f(x) = (x^2 - 4)\cdot\mathrm{e}^{-x}$ genau die Nullstellen $-2$ und $2$ hat, und berechne die Stelle des Tiefpunkts.
\teil Berechne den Tiefpunkt des Graphen von $f(x) = 0{,}25x^3 - 3x$. Gib den Hochpunkt mithilfe der Punktsymmetrie ohne Rechnung an und berechne den Abstand der beiden Punkte.
\teil Berechne die Extrempunkte des Graphen von $f(x) = -0{,}5x^3 + 1{,}5x$ und zeige, dass ihre Verbindungsgerade die Winkelhalbierende $y = x$ ist.
\end{teile}
\end{aufgabe}

\begin{aufgabe}{Ich finde den Fehler bei der Art eines Extrempunkts. Jonas bestimmt die Extrempunkte von $f(x) = x^3 - 3x^2 - 9x + 2$:}
\rechnung{f'(x) &= 3x^2 - 6x - 9 = 0 && \Leftrightarrow x = -1 \text{ oder } x = 3 \\ f''(x) &= 6x - 6 \\ f''(-1) &= -12 < 0 && \Rightarrow \text{Tiefpunkt } T(-1 \mid 7) \\ f''(3) &= 12 > 0 && \Rightarrow \text{Hochpunkt } H(3 \mid -25)}
\begin{teile}
\teil Finde den Fehler, benenne ihn und korrigiere die Rechnung.
\teil Bestimme die Extrempunkte von $g(x) = x^3 + 3x^2 - 1$.
\end{teile}
\end{aufgabe}

\begin{aufgabe}{Ich kann ohne Rechnung begründen, ob an einer Stelle ein Hoch- oder Tiefpunkt liegt.}
\begin{teile}
\teil Begründe, warum aus $f'(x_0) = 0$ allein nicht folgt, dass bei $x_0$ ein Hoch- oder Tiefpunkt liegt.
\teil Begründe, warum der größte Funktionswert auf einem Intervall $[a;\,b]$ an einer Stelle liegen kann, an der die Tangente nicht waagerecht ist.
\teil Beurteile die Aussage: „Hat der Graph einer ganzrationalen Funktion dritten Grades zwei verschiedene Stellen mit waagerechter Tangente, so hat er einen Hoch- und einen Tiefpunkt.“
\teil Beurteile die Aussage: „Rechts neben einem Tiefpunkt $T(x_0 \mid y_0)$ gilt in der Nähe $f'(x) > 0$ und $f(x) < y_0$.“
\end{teile}
\end{aufgabe}

\begin{aufgabe}{Ich kann mit der Ableitung begründen, welche Hoch- und Tiefpunkte ein Graph hat.}
\begin{teile}
\teil Gegeben ist $f(x) = x^3 - 3x^2$ mit $f'(x) = 3x\,(x - 2)$. Gib die Schnittpunkte des Graphen mit den Koordinatenachsen an und begründe mit $f'$, dass der Graph im Ursprung einen Hochpunkt hat.
\teil Begründe mit a), dass der Graph die $x$-Achse im Ursprung berührt und nicht schneidet.
\teil Gib den Tiefpunkt des Graphen von $f(x) = \mathrm{e}^{x} - x$ an und weise nach, dass der Graph keine weiteren Extrempunkte hat.
\teil Begründe mit Monotonie, Grenzverhalten und Symmetrie, dass $f(x) = \mathrm{e}^{-x^2}$ den größten Funktionswert $1$ hat, und gib die Wertemenge von $f$ an.
\end{teile}
\end{aufgabe}

\begin{aufgabe}{Ich kann Gipfelpunkte in einer Sachaufgabe berechnen und vergleichen. Das Profil einer Hügelkette wird für $0 \le x \le 4$ durch $f(x) = -0{,}5x^2 + 2x + 1$ und für $4 \le x \le 7$ durch $g(x) = -0{,}25x^3 + 3{,}75x^2 - 18x + 29$ beschrieben ($x$ und $y$ in 100 m).}
\begin{teile}
\teil Berechne die beiden Gipfelpunkte.
\teil Zwischen den Gipfeln soll ein Stahlseil geradlinig gespannt werden, das höchstens 450 m lang sein darf. Prüfe, ob das möglich ist.
\end{teile}
\end{aufgabe}
